% ============================================================
% math-paper-en  MCM/ICM (COMAP) English paper template
% Compile with XeLaTeX.  Deliverable language: English only.
% Project layout, labels and gate contract follow the math-paper-en skill;
% see references/comap-format.md for the summary sheet, section whitelist,
% page contract and appendix structure.  Do not rename the style commands
% or the section titles of this template.
% ============================================================
\documentclass[12pt,letterpaper]{article}

% ---------- fonts: English only ----------
\usepackage{fontspec}
\IfFontExistsTF{Times New Roman}
  {\setmainfont{Times New Roman}\setsansfont{Arial}}
  {\setmainfont{texgyretermes-regular.otf}[BoldFont=texgyretermes-bold.otf,ItalicFont=texgyretermes-italic.otf,BoldItalicFont=texgyretermes-bolditalic.otf]%
   \setsansfont{texgyreheros-regular.otf}[BoldFont=texgyreheros-bold.otf]}
\IfFontExistsTF{Cambria Math}{\usepackage{unicode-math}\setmathfont{Cambria Math}}{\usepackage{amsmath,amssymb,bm,mathrsfs}}

% ---------- page & spacing ----------
\usepackage[top=2.4cm, bottom=2.4cm, left=2.2cm, right=2.2cm]{geometry}
\usepackage{setspace}
\onehalfspacing
\setlength{\parindent}{2em}

% ---------- header / footer ----------
\usepackage{fancyhdr}
\usepackage{lastpage}
\pagestyle{fancy}
\fancyhf{}
\setlength{\headheight}{14pt}
\fancyfoot[C]{{\fontsize{9pt}{11pt}\selectfont Page \thepage\ of \pageref{LastPage}}}
\renewcommand{\headrulewidth}{0pt}

% ---------- figures / tables / algorithms ----------
\usepackage{graphicx}
\graphicspath{{./}{./figures/}{./images/}}
\usepackage{subcaption}
\usepackage{float}
\usepackage[table,xcdraw]{xcolor}
\usepackage{booktabs}
\usepackage{longtable}
\usepackage{array}
\usepackage{multirow}
\usepackage{algorithm}
\usepackage{algpseudocode}
\renewcommand{\algorithmicrequire}{\textbf{Input:}}
\renewcommand{\algorithmicensure}{\textbf{Output:}}

% ---------- references / links / code ----------
\usepackage[hidelinks]{hyperref}
\hypersetup{colorlinks=true,linkcolor=black,citecolor=black,urlcolor=black}
\usepackage{listings}
\lstset{
    basicstyle=\footnotesize\ttfamily,
    numbers=left,
    numberstyle=\tiny,
    frame=single,
    breaklines=true,
    language=Python,
    commentstyle=\color{gray}
}

% ---------- misc ----------
\usepackage{indentfirst}
\usepackage{enumitem}
\usepackage{makecell}
\usepackage{titlesec}
\usepackage{appendix}

% ---------- style commands reused by the gate checks ----------
\newcommand{\titlefont}{\fontsize{16pt}{20pt}\selectfont\bfseries}
\newcommand{\sectiontitlefont}{\fontsize{13.5pt}{17pt}\selectfont\bfseries}
\newcommand{\levelonefont}{\fontsize{13.5pt}{17pt}\selectfont\bfseries}
\newcommand{\leveltwofont}{\fontsize{12.5pt}{16pt}\selectfont\bfseries}
\newcommand{\levelthreefont}{\fontsize{12pt}{16pt}\selectfont\bfseries}
\newcommand{\bodyfont}{\fontsize{12pt}{16pt}\selectfont}
\newcommand{\keywordfont}{\fontsize{12pt}{16pt}\selectfont\bfseries}
\newcommand{\paperstrong}[1]{\textbf{#1}}
\newcommand{\appendixtitlefont}{\fontsize{13.5pt}{17pt}\selectfont\bfseries}
\newcommand{\stepfont}[1]{{\fontsize{12pt}{16pt}\selectfont\bfseries Step~#1}}
\newcommand{\bibfont}{\fontsize{10.5pt}{14pt}\selectfont}

\renewcommand{\thesection}{\arabic{section}}
\renewcommand{\thesubsection}{\arabic{section}.\arabic{subsection}}
\renewcommand{\thesubsubsection}{\arabic{section}.\arabic{subsection}.\arabic{subsubsection}}
\titleformat{\section}{\levelonefont}{\thesection.}{0.6em}{}
\titleformat{\subsection}{\leveltwofont}{\thesubsection}{0.6em}{}
\titleformat{\subsubsection}{\levelthreefont}{\thesubsubsection}{0.6em}{}
\setcounter{tocdepth}{2}

% ============================================================
\begin{document}
\label{body:start}

% ---------- PAGE 1: SUMMARY SHEET (no author, no affiliation) ----------
\begin{center}
\vspace*{0.6cm}
{\titlefont [Paper title: replace with the competition problem title]\par}
\vspace{0.35cm}
\end{center}

\begin{center}
{\sectiontitlefont Summary}
\end{center}

\vspace{0.25cm}
\bodyfont

\indent Replace this paragraph with the restatement of what the problem asks and what the paper delivers. The summary sheet is the single most important page: the judges read it first, so it must be self-contained, quantitative and free of formulas.

\indent \paperstrong{For Problem 1}, we build \paperstrong{[named model]} and obtain \paperstrong{[key quantitative result]}. State the model, the data, the decision criterion and the number that follows from it.

\indent \paperstrong{For Problem 2}, we extend the model with \paperstrong{[technique]} and report \paperstrong{[key quantitative result]}.

\indent \paperstrong{For Problem 3}, we formulate \paperstrong{[optimization or learning formulation]} and solve it with \paperstrong{[algorithm]}; the solution achieves \paperstrong{[key quantitative result]}.

\indent \paperstrong{For Problem 4}, we validate the model on \paperstrong{[validation data]} and quantify the uncertainty with \paperstrong{[sensitivity or error analysis]}.

\indent In summary, the paper links \paperstrong{[model family]} to the decision the problem asks for, and every claim is supported by a reproducible figure, table or numerical experiment. Replace the bracketed text; keep the page free of equations.

\indent {\keywordfont \paperstrong{Keywords:}} \paperstrong{[keyword 1]}; \paperstrong{[keyword 2]}; \paperstrong{[keyword 3]}; \paperstrong{[keyword 4]}; \paperstrong{[keyword 5]}\label{abstract:end}

\newpage

% ---------- PAGE 2: TABLE OF CONTENTS (allowed by COMAP) ----------
\tableofcontents
\newpage

\section{Introduction}
Restate the problem in your own words: background, what must be decided, inputs, outputs, constraints and the questions asked. Finish with one short paragraph that states the contribution of this paper and how the sections are organised.

\section{Assumptions and Justifications}
\begin{enumerate}
    \item State each assumption and justify it with data, physics, regulation or the problem statement.
    \item Every assumption must be used later; assumptions never used in the model are deleted.
\end{enumerate}

\section{Notations}
\begin{table}[H]
\centering
\begin{tabular}{cl}
\toprule
Symbol & Meaning \\
\midrule
$D$ & water depth or other primary state variable \\
$\theta$ & opening angle or other decision parameter \\
\bottomrule
\end{tabular}
\caption{Main notations used in this paper.}
\end{table}

\section{Model Design and Solution}
\subsection{Model for Problem 1}
Give the derivation chain before the result: define the variables, state the governing relation, derive the working formula, then solve and report the number. Insert the figure of the first question with \verb|\includegraphics| and explain it in the following paragraph.

\subsection{Model for Problem 2}
Add only what the new question changes; state explicitly which part of the previous model is reused and which is modified.

\subsection{Model for Problem 3}
Present the optimisation or learning formulation, the algorithm, the stopping criterion and the achieved objective value.

\subsection{Model for Problem 4}
Present the extended model, the validation protocol and the error or uncertainty statement.

\section{Sensitivity Analysis}
Vary the dominant parameters over defensible ranges, report how the conclusion moves, and say which conclusion is robust and which is conditional.

\section{Model Evaluation}
Write the strengths and the weaknesses of the model in two clearly separated paragraphs; a weakness statement must name the regime in which the model fails.

\section{Conclusions}
Answer each question in one sentence with its key number, then state the limits of applicability.

\begin{thebibliography}{99}
\bibfont
\bibitem{ref1} Author, A. Title of the paper or data source. Journal or Institution, Year. URL or DOI.
\bibitem{ref2} Author, B. Title of the second source. Publisher, Year.
\end{thebibliography}
\label{body:end}

\newpage
\label{appendix:start}
\begin{appendices}

\section{Supporting Material Catalog}
\begin{table}[H]
\centering
\begin{tabular}{ll}
\toprule
File & Role \\
\midrule
code\_q1.py & Cleaning and first-question computation \\
code\_q4.py & Sensitivity experiment \\
results\_q1.xlsx & Numeric results of the first question \\
\bottomrule
\end{tabular}
\end{table}

\section{Code}
\begin{lstlisting}
# full reproducible code, no comments, one file per question
import numpy as np

def main():
    pass

if __name__ == "__main__":
    main()
\end{lstlisting}

\end{appendices}
\label{paper:end}

\end{document}